% Options for packages loaded elsewhere
\PassOptionsToPackage{unicode}{hyperref}
\PassOptionsToPackage{hyphens}{url}
%
\documentclass[
  fontset=fandol]{ctexart}
\usepackage{amsmath,amssymb}
\usepackage{lmodern}
\usepackage{iftex}
\ifPDFTeX
  \usepackage[T1]{fontenc}
  \usepackage[utf8]{inputenc}
  \usepackage{textcomp} % provide euro and other symbols
\else % if luatex or xetex
  \usepackage{unicode-math}
  \defaultfontfeatures{Scale=MatchLowercase}
  \defaultfontfeatures[\rmfamily]{Ligatures=TeX,Scale=1}
\fi
% Use upquote if available, for straight quotes in verbatim environments
\IfFileExists{upquote.sty}{\usepackage{upquote}}{}
\IfFileExists{microtype.sty}{% use microtype if available
  \usepackage[]{microtype}
  \UseMicrotypeSet[protrusion]{basicmath} % disable protrusion for tt fonts
}{}
\makeatletter
\@ifundefined{KOMAClassName}{% if non-KOMA class
  \IfFileExists{parskip.sty}{%
    \usepackage{parskip}
  }{% else
    \setlength{\parindent}{0pt}
    \setlength{\parskip}{6pt plus 2pt minus 1pt}}
}{% if KOMA class
  \KOMAoptions{parskip=half}}
\makeatother
\usepackage{xcolor}
\IfFileExists{xurl.sty}{\usepackage{xurl}}{} % add URL line breaks if available
\IfFileExists{bookmark.sty}{\usepackage{bookmark}}{\usepackage{hyperref}}
\hypersetup{
  hidelinks,
  pdfcreator={LaTeX via pandoc}}
\urlstyle{same} % disable monospaced font for URLs
\usepackage[margin=1.8cm]{geometry}
\usepackage{color}
\usepackage{fancyvrb}
\newcommand{\VerbBar}{|}
\newcommand{\VERB}{\Verb[commandchars=\\\{\}]}
\DefineVerbatimEnvironment{Highlighting}{Verbatim}{commandchars=\\\{\}}
% Add ',fontsize=\small' for more characters per line
\usepackage{framed}
\definecolor{shadecolor}{RGB}{248,248,248}
\newenvironment{Shaded}{\begin{snugshade}}{\end{snugshade}}
\newcommand{\AlertTok}[1]{\textcolor[rgb]{0.94,0.16,0.16}{#1}}
\newcommand{\AnnotationTok}[1]{\textcolor[rgb]{0.56,0.35,0.01}{\textbf{\textit{#1}}}}
\newcommand{\AttributeTok}[1]{\textcolor[rgb]{0.77,0.63,0.00}{#1}}
\newcommand{\BaseNTok}[1]{\textcolor[rgb]{0.00,0.00,0.81}{#1}}
\newcommand{\BuiltInTok}[1]{#1}
\newcommand{\CharTok}[1]{\textcolor[rgb]{0.31,0.60,0.02}{#1}}
\newcommand{\CommentTok}[1]{\textcolor[rgb]{0.56,0.35,0.01}{\textit{#1}}}
\newcommand{\CommentVarTok}[1]{\textcolor[rgb]{0.56,0.35,0.01}{\textbf{\textit{#1}}}}
\newcommand{\ConstantTok}[1]{\textcolor[rgb]{0.00,0.00,0.00}{#1}}
\newcommand{\ControlFlowTok}[1]{\textcolor[rgb]{0.13,0.29,0.53}{\textbf{#1}}}
\newcommand{\DataTypeTok}[1]{\textcolor[rgb]{0.13,0.29,0.53}{#1}}
\newcommand{\DecValTok}[1]{\textcolor[rgb]{0.00,0.00,0.81}{#1}}
\newcommand{\DocumentationTok}[1]{\textcolor[rgb]{0.56,0.35,0.01}{\textbf{\textit{#1}}}}
\newcommand{\ErrorTok}[1]{\textcolor[rgb]{0.64,0.00,0.00}{\textbf{#1}}}
\newcommand{\ExtensionTok}[1]{#1}
\newcommand{\FloatTok}[1]{\textcolor[rgb]{0.00,0.00,0.81}{#1}}
\newcommand{\FunctionTok}[1]{\textcolor[rgb]{0.00,0.00,0.00}{#1}}
\newcommand{\ImportTok}[1]{#1}
\newcommand{\InformationTok}[1]{\textcolor[rgb]{0.56,0.35,0.01}{\textbf{\textit{#1}}}}
\newcommand{\KeywordTok}[1]{\textcolor[rgb]{0.13,0.29,0.53}{\textbf{#1}}}
\newcommand{\NormalTok}[1]{#1}
\newcommand{\OperatorTok}[1]{\textcolor[rgb]{0.81,0.36,0.00}{\textbf{#1}}}
\newcommand{\OtherTok}[1]{\textcolor[rgb]{0.56,0.35,0.01}{#1}}
\newcommand{\PreprocessorTok}[1]{\textcolor[rgb]{0.56,0.35,0.01}{\textit{#1}}}
\newcommand{\RegionMarkerTok}[1]{#1}
\newcommand{\SpecialCharTok}[1]{\textcolor[rgb]{0.00,0.00,0.00}{#1}}
\newcommand{\SpecialStringTok}[1]{\textcolor[rgb]{0.31,0.60,0.02}{#1}}
\newcommand{\StringTok}[1]{\textcolor[rgb]{0.31,0.60,0.02}{#1}}
\newcommand{\VariableTok}[1]{\textcolor[rgb]{0.00,0.00,0.00}{#1}}
\newcommand{\VerbatimStringTok}[1]{\textcolor[rgb]{0.31,0.60,0.02}{#1}}
\newcommand{\WarningTok}[1]{\textcolor[rgb]{0.56,0.35,0.01}{\textbf{\textit{#1}}}}
\usepackage{longtable,booktabs,array}
\usepackage{calc} % for calculating minipage widths
% Correct order of tables after \paragraph or \subparagraph
\usepackage{etoolbox}
\makeatletter
\patchcmd\longtable{\par}{\if@noskipsec\mbox{}\fi\par}{}{}
\makeatother
% Allow footnotes in longtable head/foot
\IfFileExists{footnotehyper.sty}{\usepackage{footnotehyper}}{\usepackage{footnote}}
\makesavenoteenv{longtable}
\setlength{\emergencystretch}{3em} % prevent overfull lines
\providecommand{\tightlist}{%
  \setlength{\itemsep}{0pt}\setlength{\parskip}{0pt}}
\setcounter{secnumdepth}{-\maxdimen} % remove section numbering
% --- custom header for final_report ---
% Map math/unicode symbols that the Latin roman font lacks to math-mode glyphs
\usepackage{newunicodechar}
\newunicodechar{≠}{\ensuremath{\neq}}
\newunicodechar{≈}{\ensuremath{\approx}}
\newunicodechar{≥}{\ensuremath{\geq}}
\newunicodechar{≤}{\ensuremath{\leq}}
\newunicodechar{×}{\ensuremath{\times}}
\newunicodechar{→}{\ensuremath{\rightarrow}}
\newunicodechar{←}{\ensuremath{\leftarrow}}
\newunicodechar{·}{\ensuremath{\cdot}}
\newunicodechar{…}{\ldots}

% Make every (long)table render in a smaller font so wide matrices fit
\usepackage{etoolbox}
\AtBeginEnvironment{longtable}{\footnotesize}
\AtBeginEnvironment{tabular}{\footnotesize}

% Allow page breaks inside code blocks / shaded verbatim
\usepackage{xpatch}

% Help TeX wrap long inline code / URLs without running off the page
\setlength{\emergencystretch}{3em}
\sloppy
\ifLuaTeX
  \usepackage{selnolig}  % disable illegal ligatures
\fi

\author{}
\date{}

\begin{document}

\hypertarget{ux57faux4e8eux6df1ux5ea6ux5b66ux4e60ux7684-a-ux80a1ux6b21ux65e5ux6536ux76caux9884ux6d4bux4e0eux6a21ux62dfux4ea4ux6613}{%
\section{基于深度学习的 A
股次日收益预测与模拟交易}\label{ux57faux4e8eux6df1ux5ea6ux5b66ux4e60ux7684-a-ux80a1ux6b21ux65e5ux6536ux76caux9884ux6d4bux4e0eux6a21ux62dfux4ea4ux6613}}

\begin{quote}
USTC 深度学习基础 2026 大作业 ·
小组：{[}姓名1 / 学号1{]} · {[}姓名2 /
学号2{]}
\end{quote}

\hypertarget{ux5f15ux8a00}{%
\subsection{1.
引言}\label{ux5f15ux8a00}}

金融时间序列高噪声、非平稳、强横截面依赖。本项目以「日频
A 股 Top-N
选股」为载体，构建端到端深度学习流水线：\textbf{数据整理
→ 因子 → 滑窗样本 → 深度模型 → IC 评估 →
回测 → 盘前信号 CLI}，并在 2026-06-01
\textasciitilde{} 06-12
同花顺模拟交易赛中实盘演练。

\hypertarget{ux6838ux5fc3ux8bbeux8ba1ux9009ux62e9}{%
\subsubsection{1.1
核心设计选择}\label{ux6838ux5fc3ux8bbeux8ba1ux9009ux62e9}}

\begin{longtable}[]{@{}
  >{\raggedright\arraybackslash}p{(\columnwidth - 4\tabcolsep) * \real{0.33}}
  >{\raggedright\arraybackslash}p{(\columnwidth - 4\tabcolsep) * \real{0.33}}
  >{\raggedright\arraybackslash}p{(\columnwidth - 4\tabcolsep) * \real{0.33}}@{}}
\toprule
设计点 & 选择 & 理由 \\
\midrule
\endhead
股票池 & 全 A 股去 ST +
北交所（\textasciitilde5000） &
与比赛口径一致，最大化训练样本 \\
标签 &
\texttt{log(close\_\{t+1\}/close\_t)} &
次日收益（选股目标直接相关） \\
滑窗 & T=20 &
兼顾短期技术信号与训练效率 \\
标准化 & \textbf{窗口内}
z-score（只用窗口内统计量） & 严防全量
mean/std 带入未来信息 \\
损失 & \textbf{IC
loss}（\texttt{1\ -\ pearson(pred,\ y)}，日批）
vs MSE & 前者直接优化选股排序 \\
策略 & n=10 初始等权，每日换 k=2 &
作业示例，换手率适中，便于复现 \\
\bottomrule
\end{longtable}

\hypertarget{ux4e0e-awesome-trading-agents-ux7684ux53d6ux820d}{%
\subsubsection{1.2 与
awesome-trading-agents
的取舍}\label{ux4e0e-awesome-trading-agents-ux7684ux53d6ux820d}}

调研 2025 年 LLMQuant 社区发布的
\href{https://github.com/LLMQuant/awesome-trading-agents}{Awesome
Trading Agents v0.1} 后，做出以下取舍：

\begin{itemize}
\tightlist
\item
  收录项目中 90\% 为 \textbf{LLM Agent
  辩论 +
  美股/加密/期货}（TradingAgents、AI
  Hedge Fund、AutoHedge
  等），与本作业「深度学习为核心决策」要求不兼容，仅作为参考文献。
\item
  真正可借鉴的 A
  股日频子集：\texttt{hsliuping/TradingAgents-CN}、\texttt{KylinMountain/TradingAgents-AShare}（仅抄去
  ST/北交所的清洗逻辑）、\texttt{ZhuLinsen/daily\_stock\_analysis}（IC/ICIR
  评估套路）、\texttt{Miasyster/QuantGPT}（因子假设→回测
  pipeline）。
\item
  亮点：我们编写了
  \texttt{skills/a-share-daily-report/SKILL.md}，让
  Claude Code
  一条命令跑完整日盘前流程，这是
  Anthropic 2025 新推 Agent Skills
  标准的一种本地化实践。
\end{itemize}

\hypertarget{ux6570ux636eux5904ux7406ux4e0eux95eeux9898ux5b9aux4e49}{%
\subsection{2.
数据处理与问题定义}\label{ux6570ux636eux5904ux7406ux4e0eux95eeux9898ux5b9aux4e49}}

\hypertarget{ux6570ux636eux6e90}{%
\subsubsection{2.1
数据源}\label{ux6570ux636eux6e90}}

科大云盘同步的
\texttt{documents-export-YYYY-MM-DD.zip}
解压至 \texttt{data/}：

\begin{longtable}[]{@{}
  >{\raggedright\arraybackslash}p{(\columnwidth - 4\tabcolsep) * \real{0.33}}
  >{\raggedright\arraybackslash}p{(\columnwidth - 4\tabcolsep) * \real{0.33}}
  >{\raggedright\arraybackslash}p{(\columnwidth - 4\tabcolsep) * \real{0.33}}@{}}
\toprule
目录 & 说明 & 使用 \\
\midrule
\endhead
\texttt{basic.csv} & 全部股票 meta &
过滤北交所 + 上市日 \\
\texttt{trade\_cal.csv} & 交易日历 &
日期对齐 \\
\texttt{daily/} & 日频 OHLCV + vwap &
原始量价 \\
\texttt{stock\_st/} & 每日 ST 名单 &
逐日剔除 \\
\texttt{metric/} &
基本面：PE/PB/换手率/市值等 &
进阶特征 \\
\texttt{market/} &
指数（000001.SH/000300.SH/399006.SZ） &
回测基准 \\
\bottomrule
\end{longtable}

\begin{quote}
接入策略：\texttt{src/features.py}
主线只用 OHLCV +
metric；\texttt{moneyflow/}
在两条独立支线里使用------短期
LSTM（\texttt{scripts/short\_term\_competition\_train.py}
的
\texttt{MONEYFLOW\_FEATURES}）和树模型
selected 特征集（在 \texttt{mf\_*} 衍生
rank 列上）。\texttt{news/}
文本数据未接入（工程成本高，且 Critic
阶段判断对短期选股贡献有限）。
\end{quote}

\hypertarget{ux80a1ux7968ux6c60ux8fc7ux6ee4ux4ee3ux7801srcdata_loader.pybuild_panel}{%
\subsubsection{\texorpdfstring{2.2
股票池过滤（代码：\texttt{src/data\_loader.py::build\_panel}）}{2.2 股票池过滤（代码：src/data\_loader.py::build\_panel）}}\label{ux80a1ux7968ux6c60ux8fc7ux6ee4ux4ee3ux7801srcdata_loader.pybuild_panel}}

\begin{enumerate}
\def\labelenumi{\arabic{enumi}.}
\tightlist
\item
  \texttt{basic.market\ !=\ "北交所"}
  ------ 与同花顺比赛口径一致。
\item
  逐交易日剔除当日
  \texttt{stock\_st/YYYYMMDD.csv} 中的
  ts\_code。
\item
  剔除上市 \textless{} 60
  日的新股（规避次新股炒作的异常
  pct\_chg）。
\item
  按 \texttt{trade\_cal.csv}
  只保留上交所交易日。
\end{enumerate}

\hypertarget{ux7279ux5f81ux4ee3ux7801srcfeatures.py}{%
\subsubsection{\texorpdfstring{2.3
特征（代码：\texttt{src/features.py}）}{2.3 特征（代码：src/features.py）}}\label{ux7279ux5f81ux4ee3ux7801srcfeatures.py}}

\begin{itemize}
\tightlist
\item
  \textbf{时序}（每股计算，T-1
  及之前）：ret\_1/5/20、ma\_5/20、std\_5/20、vol\_ratio\_5、amihud\_20、vwap
  偏离、RSI14、MACD 及其信号/柱。
\item
  \textbf{基本面}（metric，ffill
  规避发布滞后）：turnover\_rate、log(circ\_mv)、1/pe\_ttm、1/pb。
\item
  \textbf{横截面 rank}（日度
  {[}0,1{]}）：rk\_ret\_1、rk\_ret\_5、rk\_turn、rk\_mv。
\end{itemize}

\textbf{防泄露关键一步}：\texttt{\_per\_stock}
对每只股票计算完所有因子后统一
\texttt{.shift(1)}，保证第 t
行的特征只引用 t-1 及之前。

\hypertarget{ux6807ux7b7eux4ee3ux7801srclabels.py}{%
\subsubsection{\texorpdfstring{2.4
标签（代码：\texttt{src/labels.py}）}{2.4 标签（代码：src/labels.py）}}\label{ux6807ux7b7eux4ee3ux7801srclabels.py}}

\begin{itemize}
\tightlist
\item
  \texttt{y\_t\ =\ log(close\_\{t+1\}\ /\ close\_t)}。
\item
  剔除：t+1 停牌（close 缺失）、t+1
  估计涨停（\textbar pct\_chg\textbar{}
  ≥ 9.5\%）、t 当日指标缺失。
\item
  日度 0.5\% 两侧 winsorise 缓解极端值。
\end{itemize}

\hypertarget{ux6837ux672cux6784ux9020ux4ee3ux7801srcdataset.pywindowdataset}{%
\subsubsection{\texorpdfstring{2.5
样本构造（代码：\texttt{src/dataset.py::WindowDataset}）}{2.5 样本构造（代码：src/dataset.py::WindowDataset）}}\label{ux6837ux672cux6784ux9020ux4ee3ux7801srcdataset.pywindowdataset}}

\begin{itemize}
\tightlist
\item
  滑动窗口 T=20，每条样本是
  \texttt{(X:\ (20,\ F),\ y:\ scalar)}。
\item
  \textbf{窗口内} z-score（只对非 rank
  列）：\texttt{(sub\ -\ sub.mean(0))\ /\ sub.std(0)}，float64
  计算规避 float32 精度漂移（见
  \texttt{tests/test\_no\_leakage.py::test\_window\_zscore\_uses\_window\_only}）。
\item
  \texttt{DayBatchSampler} 每个 batch =
  一个交易日的全部样本，天然适配 IC loss
  与截面回测。
\end{itemize}

\hypertarget{ux6570ux636eux96c6ux5212ux5206}{%
\subsubsection{2.6
数据集划分}\label{ux6570ux636eux96c6ux5212ux5206}}

\begin{longtable}[]{@{}
  >{\raggedright\arraybackslash}p{(\columnwidth - 4\tabcolsep) * \real{0.33}}
  >{\raggedright\arraybackslash}p{(\columnwidth - 4\tabcolsep) * \real{0.33}}
  >{\raggedright\arraybackslash}p{(\columnwidth - 4\tabcolsep) * \real{0.33}}@{}}
\toprule
集合 & 时间范围 & 用途 \\
\midrule
\endhead
Train & 2019-01-01 \textasciitilde{}
2024-12-31（\textasciitilde6 年） &
拟合 \\
Val & 2025-01-01 \textasciitilde{}
2026-04-30（\textasciitilde16 月） &
选模型 + 回测 \\
比赛 & 2026-06-01 \textasciitilde{}
2026-06-12（10 日） & 事实上的 OOS \\
\bottomrule
\end{longtable}

\hypertarget{ux65e0ux6cc4ux9732ux9a8cux8bc1}{%
\subsubsection{2.7
无泄露验证}\label{ux65e0ux6cc4ux9732ux9a8cux8bc1}}

\texttt{tests/test\_no\_leakage.py}
三项：

\begin{enumerate}
\def\labelenumi{\arabic{enumi}.}
\tightlist
\item
  \textbf{\texttt{test\_features\_are\_strictly\_causal}}：在未来某天
  poison
  \texttt{close\ =\ 1e6}，验证该天特征值不爆炸（因已
  shift）。
\item
  \textbf{\texttt{test\_label\_is\_next\_day\_return}}：断言
  \texttt{y\ =\ log(close\_\{t+1\}/close\_t)}。
\item
  \textbf{\texttt{test\_window\_zscore\_uses\_window\_only}}：断言每个窗口
  z-score 均值 \textless{} 1e-5。
\end{enumerate}

CI：\texttt{pytest\ tests/} 通过。

\hypertarget{ux6a21ux578bux5b9eux73b0}{%
\subsection{3.
模型实现}\label{ux6a21ux578bux5b9eux73b0}}

我们一开始严格按作业示例从神经网络做起，过程中发现
LSTM/Transformer 在 2026Q2
严重失稳后，又补充了 LightGBM/XGBoost
一条主线。两条线代码各自独立，便于横向对比。

\hypertarget{ux6df1ux5ea6ux5b66ux4e60ux6a21ux578bux4ee3ux7801srcmodels}{%
\subsubsection{\texorpdfstring{3.1
深度学习模型（代码：\texttt{src/models/}）}{3.1 深度学习模型（代码：src/models/）}}\label{ux6df1ux5ea6ux5b66ux4e60ux6a21ux578bux4ee3ux7801srcmodels}}

三个模型按复杂度递增，共享输入/输出接口（\texttt{(B,\ T,\ F)\ →\ (B,)}）：

\begin{longtable}[]{@{}
  >{\raggedright\arraybackslash}p{(\columnwidth - 4\tabcolsep) * \real{0.33}}
  >{\raggedright\arraybackslash}p{(\columnwidth - 4\tabcolsep) * \real{0.33}}
  >{\raggedright\arraybackslash}p{(\columnwidth - 4\tabcolsep) * \real{0.33}}@{}}
\toprule
模型 & 参数规模 & 描述 \\
\midrule
\endhead
\texttt{MLP} & \textasciitilde0.1 M &
基线：flatten(T·F) → 128·3 层 → 1 \\
\texttt{LSTMModel} & 0.1
\textasciitilde{} 1.5 M & 2-3 层 LSTM,
hidden 64/128/256/384，取最后一步 \\
\texttt{TransformerModel} & 0.1
\textasciitilde{} 3 M & 2-4 层 Encoder,
d=64/128/256, 4-8 头，学习 {[}CLS{]} +
可学习位置编码 \\
\bottomrule
\end{longtable}

\hypertarget{ux635fux5931ux51fdux6570}{%
\paragraph{3.1.1
损失函数}\label{ux635fux5931ux51fdux6570}}

\begin{itemize}
\tightlist
\item
  \texttt{MSE}：标准回归。
\item
  \textbf{\texttt{IC\ loss}}：\texttt{1\ -\ pearson(pred,\ y)}，batch
  =
  一个交易日的整截面。直接优化排序而非绝对收益，更适合选股。
\item
  短期模型还引入了\textbf{指数 recency
  权重}（\texttt{scripts/short\_term\_competition\_train.py}）：训练日
  IC loss 按
  \texttt{exp(-ln2\ ·\ age\ /\ half\_life)}
  缩放，半衰期 30/60/90 天。
\end{itemize}

\hypertarget{ux8badux7ec3ux4ee3ux7801srctrain.py}{%
\paragraph{\texorpdfstring{3.1.2
训练（代码：\texttt{src/train.py}）}{3.1.2 训练（代码：src/train.py）}}\label{ux8badux7ec3ux4ee3ux7801srctrain.py}}

\begin{itemize}
\tightlist
\item
  优化器
  AdamW，lr=1e-3，wd=1e-5，梯度裁剪
  1.0。
\item
  8 个 epoch，按验证 IC 保存 best。
\item
  A100 主实验用
  \texttt{-\/-batch-cap\ 8192\ -\/-amp}，单截面单批，DayBatchSampler
  保证「一个 batch = 一个交易日」。
\item
  所有
  checkpoint、验证预测、训练日志写到
  \texttt{checkpoints/}。
\end{itemize}

\hypertarget{ux68afux5ea6ux63d0ux5347ux6811ux4e3bux7ebfux4ee3ux7801scriptsexplore_models.pyscriptstrain_lgbm_wq_short.py}{%
\subsubsection{\texorpdfstring{3.2
梯度提升树主线（代码：\texttt{scripts/explore\_models.py}、\texttt{scripts/train\_lgbm\_wq\_short.py}）}{3.2 梯度提升树主线（代码：scripts/explore\_models.py、scripts/train\_lgbm\_wq\_short.py）}}\label{ux68afux5ea6ux63d0ux5347ux6811ux4e3bux7ebfux4ee3ux7801scriptsexplore_models.pyscriptstrain_lgbm_wq_short.py}}

\begin{itemize}
\tightlist
\item
  \textbf{特征侧}：在 §2.3 时序 + 基本面
  + 横截面 rank 之外，构造了
  \textasciitilde40 个 WorldQuant 风格
  alpha（截面
  rank、量价回归、流动性、换手率），以及四个特征集合
  \texttt{base\ /\ alpha\_only\ /\ base+alpha\ /\ selected}（\texttt{selected}
  是按训练期 RankIC
  单调筛掉低贡献后留下的 80 列）。
\item
  \textbf{目标侧}：

  \begin{itemize}
  \tightlist
  \item
    LightGBM
    \textbf{LambdaRank}：每个交易日构造一个
    group，标签由次日收益分位映射成 0-N
    整数。
  \item
    LightGBM \textbf{huber} / XGBoost
    \textbf{pseudoHuber}：回归形式，方便和
    IC loss 做对照。
  \end{itemize}
\item
  \textbf{样本权重}：与深度模型对齐，按训练期日期做指数
  recency
  衰减，半衰期参数化为搜索维度（10 / 20
  / 30 / 45 / 60 / 90）。
\item
  \textbf{方向性}：每个候选额外有
  \texttt{forward\ /\ inverse}
  两个版本，用于发现「反向短期
  alpha」是否存在；最终我们用此排除了一批形似过拟合的反向候选。
\end{itemize}

\hypertarget{ux6211ux4eecux5c1dux8bd5ux8fc7ux7684ux6a21ux578bux4e0eux8badux7ec3ux65b9ux6848}{%
\subsection{4.
我们尝试过的模型与训练方案}\label{ux6211ux4eecux5c1dux8bd5ux8fc7ux7684ux6a21ux578bux4e0eux8badux7ec3ux65b9ux6848}}

作业要求重点写「实验亮点」。本节先用一张总表把所有跑过的方案汇总（尝试
→ 目的 → 最好结果 → 是否进入最终方案 →
淘汰原因），再展开三条主线：\textbf{4.1
深度学习失败 / 4.2 树模型胜出 / 4.3
alpha-research 补充未
promote}。所有数值都可在
\texttt{reports/compare\_metrics.csv}、\texttt{reports/may\_2026\_validation/summary\_*.csv}、\texttt{alpha-stage/artifacts/alpha\_results.json}
复算。

\hypertarget{ux603bux8868ux5c1dux8bd5ux77e9ux9635}{%
\subsubsection{4.0
总表：尝试矩阵}\label{ux603bux8868ux5c1dux8bd5ux77e9ux9635}}

\begin{longtable}[]{@{}
  >{\raggedright\arraybackslash}p{(\columnwidth - 10\tabcolsep) * \real{0.17}}
  >{\raggedright\arraybackslash}p{(\columnwidth - 10\tabcolsep) * \real{0.17}}
  >{\raggedright\arraybackslash}p{(\columnwidth - 10\tabcolsep) * \real{0.17}}
  >{\raggedright\arraybackslash}p{(\columnwidth - 10\tabcolsep) * \real{0.17}}
  >{\raggedright\arraybackslash}p{(\columnwidth - 10\tabcolsep) * \real{0.17}}
  >{\raggedright\arraybackslash}p{(\columnwidth - 10\tabcolsep) * \real{0.17}}@{}}
\toprule
\# & 尝试 & 目的 & 最好结果 &
入选最终方案 & 淘汰原因 \\
\midrule
\endhead
A & MLP + MSE & 作业示例基线 & 长验证 IC
0.014, Top10 +13 bp & 否 &
信号近零，作为 baseline 比较 \\
B & Transformer + IC loss &
看注意力是否更适合横截面 & 长验证 IC
0.005\textasciitilde0.020 & 否 & 4 层
d=256 完全失效，2 层 mid 也仅 0.020 \\
C & LSTM + IC loss + DayBatch & 作业主线
& 长验证 IC \textbf{0.099} / RankIC
0.090 / Top10 +173 bp & 否 & 5 月 OOS
年化 -88.9\%，regime shift 后反向 \\
D & 短期 LSTM + recency 衰减 + 资金流 &
用近端样本对齐 5 月 & 4 个短期 LSTM May
年化均 \textless{} -70\% & 否 &
短数据无法救 LSTM，整条 LSTM 路被否 \\
E & sklearn SGD（线性 ranker baseline）
& 给树模型一个最简单参考线 & RankIC
0.049（一度全场最高） & 否 & May BT 年化
-47.4\%，RankIC ≠ 收益 \\
F & LightGBM \texttt{lgbm\_wq\_*} 七候选
& 把信号搬到树模型 + WQ alpha &
wq\_06（huber+hl30）May 年化 +12.1\% &
否（被 4.2.2 的网格替代） &
是阶段性最优，后被 explore\_002
全面超越 \\
G & 43-候选探索网格（LGBM/XGB × 损失 ×
特征 × HL × 起点） & 系统化扫调参空间 &
explore\_002（LambdaRank+selected+hl30+train2025）BT
Sharpe \textbf{2.38}, BT 年化
\textbf{+144\%} &
\textbf{是}，作为最终上线模型 & --- \\
H & fast 因子分组
LGBM（按论文/语义分组） &
看是否单组就够强 &
\texttt{wq\_momentum\ +\ huber\ +\ forward}
Sharpe 2.50 / +162\% & 否 & 仅 15 个 BT
日，过短；与 explore\_002 高度同源 \\
I & 文献组合（Carhart, FF+, AQR, MSCI
Barra, Qlib158）huber &
文献已知风格因子能否在 A 股 5 月可用 &
\texttt{worldquant\_101\_price\_volume}（huber+forward）Sharpe
1.05 /
+37\%；\texttt{carhart\_ff\_plus\_momentum}（huber+forward）IC
0.044 / Sharpe -0.13 / -21\% & 否 & 多数
Sharpe 在 0\textasciitilde1 之间，不优于
explore\_002；\texttt{carhart} IC
正但回测仍亏，再次说明信号 ≠ 收益 \\
J & rescue\_20260602（5-6 天近窗 tune）
& 6/2 现场救火，挑能交易的 6/3 名单 &
\texttt{rescue\_2025\_rank\_alpha} 6 天
IC 0.047 / Sharpe 8.3 & 否 & 6 个 BT
日太短；同表 RankIC 反向（-0.013），同时
huber 版 RankIC 也是负的 \\
K & event reversal
模型（恐慌反弹分类器） & 把次日反弹作为
0/1 分类，避开回归 & val AUC
\textbf{0.626} & 否 & 但 top-k 5/10/20
名单的 mean\_next\_pct
全部为负（-0.46\textasciitilde-0.55），AUC
高 ≠ 选股可用 \\
L & LightGBM \texttt{lr\_*}
学习率/树数网格（24 候选 ×4 cost） &
检查是否有低 lr / 多树的稳定收益 &
全部条目均落后于 G & 否 & 没人胜出 \\
M & \texttt{alpha-stage} Cycle
1（A001\textasciitilde A006） &
显式因子库 + agent pipeline & A003 20D
动量短窗 RankIC 0.10、A005 Amihud RankIC
0.029（短窗假阳性） & 否 & 扩窗后 RankIC
反向，Codex 二审 1/10，A003/A005 kill \\
N & \texttt{alpha-stage} Cycle
2（A019\textasciitilde A030
中期动量+低波动+流动性等） &
修一遍后再选 repair 候选 & A029
中期动量短回撤低冲击：H5 train RankIC
0.069 / val RankIC \textbf{0.045} / 长短
5bps 4.9\% / 30bps -2.2\%，long-only
净收益 14.4\%\textasciitilde18.1\% &
否（保留为 repair） & 30bps
长短转负，未过 critic gate；只作 spread
诊断 \\
O & 跨模型 ensemble（explore\_002 + 004
+ 029 cross-section rank 平均） &
计划用差异化打分降方差 & 仅 5/29
离线生成过
\texttt{20260529\_ensemble\_targets.csv}
& \textbf{没真正部署到 6/1} & 6/1
凌晨只有 explore\_002 训练完，004/029 因
GPU 排队/数据 schema
临时回退；最终上线只用 explore\_002 \\
\bottomrule
\end{longtable}

总表读法：A\textasciitilde D 是「为什么
DL 不行」，E\textasciitilde L
是「为什么树模型这条线最后只剩
G」，M\textasciitilde N 是「另一条独立
alpha 实验线，没产出 promote
但产出了边界」，O 是「计划中的 ensemble
没真正打到比赛」。下面三个小节按这三条主线展开。

\hypertarget{ux4e3bux7ebf-1ux6df1ux5ea6ux5b66ux4e60ux8defux7ebfux88ab-regime-shift-ux5426ux51b3}{%
\subsubsection{4.1 主线
1：深度学习路线被 regime shift
否决}\label{ux4e3bux7ebf-1ux6df1ux5ea6ux5b66ux4e60ux8defux7ebfux88ab-regime-shift-ux5426ux51b3}}

\hypertarget{ux957fux9a8cux8bc1ux96c6ux4e0aux7684ux6f02ux4eaeux6570ux5b57}{%
\paragraph{4.1.1
长验证集上的「漂亮数字」}\label{ux957fux9a8cux8bc1ux96c6ux4e0aux7684ux6f02ux4eaeux6570ux5b57}}

按作业示例先训练 MLP / LSTM /
Transformer 三个共享接口的小模型（输入
\texttt{(B,\ 20,\ F)} → 输出
\texttt{(B,)}），损失对比 MSE 与 IC
loss，验证窗 2025-01 \textasciitilde{}
2026-04（\textasciitilde300
个交易日，记作 long val）：

\begin{longtable}[]{@{}
  >{\raggedright\arraybackslash}p{(\columnwidth - 10\tabcolsep) * \real{0.17}}
  >{\raggedright\arraybackslash}p{(\columnwidth - 10\tabcolsep) * \real{0.17}}
  >{\raggedright\arraybackslash}p{(\columnwidth - 10\tabcolsep) * \real{0.17}}
  >{\raggedright\arraybackslash}p{(\columnwidth - 10\tabcolsep) * \real{0.17}}
  >{\raggedright\arraybackslash}p{(\columnwidth - 10\tabcolsep) * \real{0.17}}
  >{\raggedright\arraybackslash}p{(\columnwidth - 10\tabcolsep) * \real{0.17}}@{}}
\toprule
Tag & IC & RankIC & ICIR & Top-10 利差
(bp) & 备注 \\
\midrule
\endhead
\texttt{mlp\_mse} & 0.014 & 0.023 & 4.01
& +13 & MSE 基线，几乎无信号 \\
\texttt{transformer\_ic\_large} & 0.005
& 0.001 & 2.32 & +7 & 4 层 d=256，IC
loss，几乎全噪声 \\
\texttt{transformer\_ic\_mid} & 0.020 &
0.014 & 8.97 & -22 & 中等
Transformer，long-only 利差转负 \\
\texttt{lstm\_ic\_large} &
\textbf{0.099} & 0.088 & 13.28 & +169 &
2 层 hidden=256，IC loss \\
\texttt{lstm\_ic\_h384\_l2\_w20} & 0.099
& \textbf{0.090} & 13.24 & \textbf{+173}
& hidden=384，RankIC/利差最佳 \\
\texttt{lstm\_ic\_h384\_l3\_w20} & 0.098
& 0.089 & \textbf{13.82} & +145 & 3
层，ICIR 最好 \\
\texttt{lstm\_ic\_h256\_l3\_w40} & 0.092
& 0.083 & 13.58 & +152 & window=40 \\
\bottomrule
\end{longtable}

\textbf{深度学习这条线内部的小结论}（这一段只在
LSTM/MLP/Transformer
之间作对比，不外推到树模型）：

\begin{enumerate}
\def\labelenumi{\arabic{enumi}.}
\tightlist
\item
  IC loss + DayBatchSampler 显著好于
  MSE：MLP+MSE IC≈0.014，换成 LSTM+IC
  loss 直接到
  0.099，与「目标函数和评估指标对齐」的直觉一致；后续所有深度模型默认走这条路。
\item
  大 Transformer 不工作：4 层 256d 在
  \textasciitilde6 年训练数据上
  IC≈0.005，2 层中等版本也只有
  0.020。结合 SNR 极低 +
  自注意力对噪声敏感，我们把 Transformer
  列为「这套数据上无法训出来」。
\item
  LSTM hidden 从 256 升到 384、层数从 2
  升到 3 在 long val 上几乎不再提升
  IC，但 RankIC、ICIR 略升，让我们暂时把
  LSTM hidden=384/2 层当成主线候选。
\end{enumerate}

\hypertarget{ux6708ux5355ux6708ux9a8cux8bc1ux6240ux6709-lstm-ux5168ux70b8}{%
\paragraph{4.1.2 5 月单月验证：所有 LSTM
全炸}\label{ux6708ux5355ux6708ux9a8cux8bc1ux6240ux6709-lstm-ux5168ux70b8}}

把上一节的 LSTM 拿到 2026 年 5 月（16
个交易日，short val）单独跑回测：

\begin{longtable}[]{@{}
  >{\raggedright\arraybackslash}p{(\columnwidth - 8\tabcolsep) * \real{0.20}}
  >{\raggedright\arraybackslash}p{(\columnwidth - 8\tabcolsep) * \real{0.20}}
  >{\raggedright\arraybackslash}p{(\columnwidth - 8\tabcolsep) * \real{0.20}}
  >{\raggedright\arraybackslash}p{(\columnwidth - 8\tabcolsep) * \real{0.20}}
  >{\raggedright\arraybackslash}p{(\columnwidth - 8\tabcolsep) * \real{0.20}}@{}}
\toprule
Tag & May IC & May RankIC & May 年化 &
May Sharpe \\
\midrule
\endhead
\texttt{lstm\_ic\_large} & -0.024 &
-0.018 & -88.9\% & -6.64 \\
\texttt{lstm\_ic\_short\_2025\_w5\_h64\_l1\_hl60\_lr3e4}
& --- & --- & -78.2\% & -6.31 \\
\texttt{lstm\_ic\_short\_mf\_decay\_h256\_l2\_w10}
& --- & --- & -82.9\% & -4.36 \\
\texttt{lstm\_ic\_short\_w5\_h128\_hl90}
& --- & --- & -74.4\% & -3.48 \\
\bottomrule
\end{longtable}

所有 long val 表现好的 LSTM 在 May
上年化 \textless{}
-70\%。我们随后做了三件事来确认这不是单点失败：

\begin{itemize}
\tightlist
\item
  训练「短窗口 + recency-weighted IC
  loss」的小 LSTM（半衰期 30/60/90
  天），引入资金流特征
  \texttt{mf\_net\_amt\_ratio}：四个候选
  May 年化均 \textless{} -70\%。
\item
  把 LSTM
  的方向反转作为对照（\texttt{*\_inverse}）：表面好看，但等同于「拿一个已知反向的信号当主信号」，没有可解释性，被否。
\item
  用一个最简单的 sklearn SGD 线性
  ranker（recency 衰减、特征同
  selected）作 baseline：RankIC
  0.049（彼时全场最高），但 May BT 年化
  -47.4\% / Sharpe -2.09 --- RankIC
  高但回测亏，再次印证「short val
  上分布偏移很大」。
\end{itemize}

\textbf{结论}（限定到
LSTM/Transformer/SGD 这条线）：在 long
val 上 IC=0.099 的 LSTM，到 5 月单月就是
-88.9\% 年化。换句话说，long val 上的 IC
不能直接外推到比赛期；深度模型本身没问题，但样本分布漂移让
2019-2024 学到的因子在 2026Q2
反向，这条路不能直接用。

\hypertarget{ux4e3bux7ebf-2ux6811ux6a21ux578b-wq-ux98ceux683cux7279ux5f81ux9010ux6b65ux80dcux51fa}{%
\subsubsection{4.2 主线 2：树模型 + WQ
风格特征逐步胜出}\label{ux4e3bux7ebf-2ux6811ux6a21ux578b-wq-ux98ceux683cux7279ux5f81ux9010ux6b65ux80dcux51fa}}

\hypertarget{ux4e03ux5019ux9009-lgbmlgbm_wq_ux5148ux786eux8ba4ux65b9ux5411}{%
\paragraph{\texorpdfstring{4.2.1 七候选
LGBM（\texttt{lgbm\_wq\_*}）：先确认方向}{4.2.1 七候选 LGBM（lgbm\_wq\_*）：先确认方向}}\label{ux4e03ux5019ux9009-lgbmlgbm_wq_ux5148ux786eux8ba4ux65b9ux5411}}

把信号源从「LSTM 处理 (B, 20,
F)」换成「LightGBM 直接学 cross-section
排序」，特征加进 \textasciitilde40 个
WorldQuant 风格 alpha（截面
rank、量价回归、流动性、换手率），训练
2025-01 \textasciitilde{} 2026-04，验证
May：

\begin{longtable}[]{@{}
  >{\raggedright\arraybackslash}p{(\columnwidth - 14\tabcolsep) * \real{0.12}}
  >{\raggedright\arraybackslash}p{(\columnwidth - 14\tabcolsep) * \real{0.12}}
  >{\raggedright\arraybackslash}p{(\columnwidth - 14\tabcolsep) * \real{0.12}}
  >{\raggedright\arraybackslash}p{(\columnwidth - 14\tabcolsep) * \real{0.12}}
  >{\raggedright\arraybackslash}p{(\columnwidth - 14\tabcolsep) * \real{0.12}}
  >{\raggedright\arraybackslash}p{(\columnwidth - 14\tabcolsep) * \real{0.12}}
  >{\raggedright\arraybackslash}p{(\columnwidth - 14\tabcolsep) * \real{0.12}}
  >{\raggedright\arraybackslash}p{(\columnwidth - 14\tabcolsep) * \real{0.12}}@{}}
\toprule
Tag & Obj & 半衰期 & May IC & May RankIC
& Top10 (bp) & May 年化 & May Sharpe \\
\midrule
\endhead
\textbf{wq\_06} & huber & 30 & +0.041 &
+0.029 & +204 & +12.1\% & +0.47 \\
wq\_01 & l1 & 30 & +0.026 & +0.015 & +29
& -22.9\% & -0.66 \\
wq\_03 & l1 & 60 & +0.041 & +0.021 & -4
& -67.8\% & -3.55 \\
\bottomrule
\end{longtable}

读出来的几条\textbf{这一阶段内}的现象（不下大结论）：

\begin{itemize}
\tightlist
\item
  同样的 IC（0.041 vs 0.041），但 RankIC
  不同（0.029 vs 0.021），huber 在 May
  BT 上比 l1 多出约 80 个百分点 --- 单看
  IC 或单看 RankIC 都会误判。
\item
  半衰期 60 → 30 把 May BT 从 -67.8\%
  拉回 +12.1\%，提示在 short val
  这种近端 16
  天的窗口上，样本权重越靠近近端越对。
\item
  wq\_06
  是「树模型这条线的第一个能用配置」，但
  Sharpe 0.47
  并不强；只能作为下一步网格的起点。
\end{itemize}

\hypertarget{ux5019ux9009ux63a2ux7d22ux7f51ux683c}{%
\paragraph{4.2.2
43-候选探索网格}\label{ux5019ux9009ux63a2ux7d22ux7f51ux683c}}

为系统化排除运气成分，我们在 A800
上跑了一次 27.7
分钟的网格搜索（\texttt{scripts/explore\_models.py}，job
27750），交叉 5 个轴：模型框架（LightGBM
/ XGBoost）× 损失（LambdaRank / huber /
pseudoHuber）× 特征集（base /
alpha\_only / base+alpha / selected）×
半衰期（10 / 20 / 30 / 45 / 60 / 90
天）× 训练起点（2024 / 2025）。43
个候选全部完成，原始结果在
\texttt{reports/may\_2026\_validation/summary\_explore.csv}。

按 May 回测 Sharpe 排序，\textbf{剔除
explore\_038
这种过拟合反例}（见下方独立小节）后的
Top 6：

\begin{longtable}[]{@{}
  >{\raggedright\arraybackslash}p{(\columnwidth - 20\tabcolsep) * \real{0.09}}
  >{\raggedright\arraybackslash}p{(\columnwidth - 20\tabcolsep) * \real{0.09}}
  >{\raggedright\arraybackslash}p{(\columnwidth - 20\tabcolsep) * \real{0.09}}
  >{\raggedright\arraybackslash}p{(\columnwidth - 20\tabcolsep) * \real{0.09}}
  >{\raggedright\arraybackslash}p{(\columnwidth - 20\tabcolsep) * \real{0.09}}
  >{\raggedright\arraybackslash}p{(\columnwidth - 20\tabcolsep) * \real{0.09}}
  >{\raggedright\arraybackslash}p{(\columnwidth - 20\tabcolsep) * \real{0.09}}
  >{\raggedright\arraybackslash}p{(\columnwidth - 20\tabcolsep) * \real{0.09}}
  >{\raggedright\arraybackslash}p{(\columnwidth - 20\tabcolsep) * \real{0.09}}
  >{\raggedright\arraybackslash}p{(\columnwidth - 20\tabcolsep) * \real{0.09}}
  >{\raggedright\arraybackslash}p{(\columnwidth - 20\tabcolsep) * \real{0.09}}@{}}
\toprule
Tag & 框架 & Obj & 特征 & HL & 训练起点
& IC & RankIC & Top10 (bp) & BT Sharpe &
BT 年化 \\
\midrule
\endhead
\textbf{explore\_002} & LGBM &
LambdaRank & selected & 30 & 2025 &
+0.075 & +0.046 & +177 & \textbf{2.38} &
\textbf{+144\%} \\
\textbf{explore\_004} & LGBM &
LambdaRank & alpha\_only & 30 & 2024 &
+0.058 & +0.029 & +165 & 1.95 &
+102\% \\
\textbf{explore\_035} & LGBM & huber &
selected & 30 & 2025 & +0.037 & +0.024 &
+148 & 1.69 & +91\% \\
explore\_039 & LGBM & huber & selected &
20 & 2025 & +0.041 & +0.026 & +186 &
1.35 & +67\% \\
\textbf{explore\_029} & XGBoost &
pseudoHuber & alpha\_only & 60 & 2025 &
+0.046 & \textbf{+0.052} & +258 & 1.19 &
+49\% \\
explore\_008 & LGBM & LambdaRank &
base+alpha & 30 & 2024 & +0.048 & +0.022
& +117 & 0.76 & +25\% \\
\bottomrule
\end{longtable}

从这次网格能稳稳读出来的几条结论（每条都用网格内的同条件对比，不再外推）：

\begin{enumerate}
\def\labelenumi{\arabic{enumi}.}
\tightlist
\item
  \textbf{同特征/同半衰期/同训练起点下，LambdaRank
  \textgreater{} huber}：explore\_002
  (LambdaRank, selected, hl30, 2025) 的
  Sharpe 2.38 vs explore\_035 (huber,
  selected, hl30, 2025) 的 Sharpe
  1.69。差距来自 LambdaRank
  直接优化排序损失，目标函数与 Top-N
  选股一致。
\item
  \textbf{2024 数据是有毒的}：完全相同的
  LambdaRank+selected+hl30，train2025 →
  Sharpe \textbf{2.38}，train2024 →
  Sharpe
  \textbf{-3.06}（explore\_000）。这是网格里最干净的对照之一。我们因此把训练起点从「2019-01」收紧到「2025-01」，牺牲样本量换分布一致性。
\item
  \textbf{XGBoost 提供异源
  RankIC}：explore\_029 的 RankIC 0.052
  是 43 个候选里最高的，但 May BT 只有
  Sharpe 1.19；它和 LightGBM
  的误差结构互补，是计划中 ensemble
  的多样性来源。
\item
  \textbf{更强正则化救了 huber}：在 wq
  七候选里 wq\_06 (huber,
  reg\_alpha=0.05, num\_leaves=31)
  Sharpe 0.47，到 explore\_035 (huber,
  reg\_alpha=0.2, num\_leaves=63) Sharpe
  1.69 ---
  同损失、同特征下，正则化是关键。
\item
  \textbf{WQ alpha
  已经携带主要信号}：alpha\_only（去掉
  OHLCV 直接特征）依旧能到 Sharpe
  1.95（explore\_004），证明信号主要来自横截面
  alpha 而非原始量价。
\end{enumerate}

\hypertarget{ux8fc7ux62dfux5408ux53cdux4f8bexplore_038ux72ecux7acbux5c0fux8282ux907fux514dux8befux5bfc}{%
\paragraph{4.2.3
过拟合反例：explore\_038（独立小节，避免误导）}\label{ux8fc7ux62dfux5408ux53cdux4f8bexplore_038ux72ecux7acbux5c0fux8282ux907fux514dux8befux5bfc}}

\texttt{explore\_038}（LGBM, huber,
selected, \textbf{inverse}, hl=10,
2025）BT Sharpe 3.09 / BT 年化
+162\%，但它的 IC 是 -0.015、RankIC
-0.029、Top10 利差 -57
bp/天。读法：把一个负向信号取反、再用 10
天极短半衰期严重偏向最后两周，等于在 16
天 BT 窗口上做了过拟合，单看 Sharpe
会得到完全错误的结论。\textbf{已
blacklist 不入选
ensemble}，本节单列以警示「永远不能只看夏普」。

\hypertarget{ux5468ux8fb9ux8865ux4e01ux5b9eux9a8cfast-ux5206ux7ec4-ux6587ux732eux7ec4ux5408-ux5b66ux4e60ux7387ux7f51ux683c-rescue-event-ux5206ux7c7bux5668}{%
\paragraph{4.2.4 周边补丁实验：fast 分组
/ 文献组合 / 学习率网格 / rescue / event
分类器}\label{ux5468ux8fb9ux8865ux4e01ux5b9eux9a8cfast-ux5206ux7ec4-ux6587ux732eux7ec4ux5408-ux5b66ux4e60ux7387ux7f51ux683c-rescue-event-ux5206ux7c7bux5668}}

为避免「explore\_002
是网格运气」，我们又跑了几组独立实验作交叉印证：

\begin{itemize}
\tightlist
\item
  \textbf{fast 因子分组
  LGBM}（\texttt{summary\_factor\_lgbm\_fast.csv}）：按语义分组单独训练。最强
  \texttt{fast\_wq\_momentum\_huber\_forward}
  Sharpe \textbf{2.50} / 年化
  \textbf{+162\%} --- 数字漂亮但只有 15
  个 BT 日，且与 explore\_002 的
  selected
  特征集高度同源，无法独立支持「另一类特征也行」的结论。
\item
  \textbf{文献组合}（\texttt{summary\_factor\_lgbm\_literature\_combos\_*.csv}）：把
  Carhart 4 因子 + 动量、Fama-French +
  动量、AQR Style、MSCI Barra Core
  Style、Qlib Alpha158-like 各打成一组用
  huber 学。最好的是
  \texttt{worldquant\_101\_price\_volume\_huber\_forward}（IC
  0.048, RankIC 0.039, Sharpe 1.05, 年化
  +37\%）；\texttt{carhart\_ff\_plus\_momentum\_huber\_forward}
  IC 0.044 / RankIC 0.031 但 May 年化
  -21\% / Sharpe -0.13 --- 又一例 IC
  正但回测亏的「单一指标陷阱」，且
  inverse 版本 Sharpe 跌到
  -5.1。文献组合整体不优于
  explore\_002，未入选。
\item
  \textbf{\texttt{lr\_*}
  学习率/树数网格}（24 个 LightGBM × 4
  个 cost=5/10/20/30bps，目录
  \texttt{backtest\_overnight\_lr\_*}）：调小
  lr、加多 num\_iterations
  换稳定性，没有候选超过
  explore\_002，作为反方向稳健性检查。
\item
  \textbf{\texttt{rescue\_20260602}}（5-6
  天近窗紧急
  tune，\texttt{reports/rescue\_20260602/}）：6/2
  凌晨想抢 6/3 的名单，用 2025-01 起 /
  2026-01 起两个起点 × rank/huber ×
  selected/alpha/base+alpha 共 9 组，仅
  6 个 BT 日。最高
  \texttt{rescue\_2025\_rank\_alpha} IC
  0.047 / Sharpe 8.3 / 年化
  +1473\%，但同表 RankIC -0.013，huber
  版 RankIC 也都是负的 --- 6
  天窗口本身不可信，\textbf{未入选实盘}，仅留作
  rescue path 备份。
\item
  \textbf{event reversal
  分类器}（\texttt{reports/reversal\_event\_20260602/}）：把「连续大跌后是否次日反弹」做成
  0/1 分类（GBDT，27
  个量价/资金流/换手率特征），val AUC
  \textbf{0.626}；但 top-k=5/10/20
  名单的 mean\_next\_pct 分别为 -0.46\%
  / -0.48\% / -0.55\%，\textbf{AUC 高 ≠
  选股可用}。这条路没有用。
\item
  \textbf{sklearn SGD 线性
  baseline}（\texttt{summary\_sklearn.csv}）：RankIC
  \textbf{0.049}（一度全场最高）但 May
  BT 年化 -47.4\% / Sharpe -2.09 --- 与
  carhart、wq\_03、explore\_038
  共同构成「单一指标陷阱」证据链。
\end{itemize}

\hypertarget{ensemble-ux8ba1ux5212ux4e0eux6700ux7ec8ux4e0aux7ebfux73b0ux5b9e}{%
\paragraph{4.2.5 Ensemble
计划与最终上线现实}\label{ensemble-ux8ba1ux5212ux4e0eux6700ux7ec8ux4e0aux7ebfux73b0ux5b9e}}

最初计划按差异化打分做 ensemble：

\begin{longtable}[]{@{}
  >{\raggedright\arraybackslash}p{(\columnwidth - 8\tabcolsep) * \real{0.20}}
  >{\raggedright\arraybackslash}p{(\columnwidth - 8\tabcolsep) * \real{0.20}}
  >{\raggedright\arraybackslash}p{(\columnwidth - 8\tabcolsep) * \real{0.20}}
  >{\raggedright\arraybackslash}p{(\columnwidth - 8\tabcolsep) * \real{0.20}}
  >{\raggedright\arraybackslash}p{(\columnwidth - 8\tabcolsep) * \real{0.20}}@{}}
\toprule
模型 & 框架 & 损失 & 特征 & 计划角色 \\
\midrule
\endhead
\texttt{explore\_002} & LightGBM &
LambdaRank & selected (80) & 主排序器 \\
\texttt{explore\_004} & LightGBM &
LambdaRank & alpha\_only (60) & WQ alpha
专家 \\
\texttt{explore\_029} & XGBoost &
pseudoHuber & alpha\_only (60) &
异源回归打分 \\
\bottomrule
\end{longtable}

合并方式是\textbf{对每个模型预测做日度
cross-section rank，再等权平均得到组合
score}（不是直接平 logit，避免 XGBoost
数值偏离主导；注意：是「rank
等权平均」，旧稿写成「等权 RankIC
加权平均」是错的，已更正）。

5/29 用三模型生成过一次离线 ensemble
名单（\texttt{reports/daily\_logs/20260529\_ensemble\_targets.csv}），但
6/1 凌晨只有 explore\_002 完成训练并落到
\texttt{checkpoints/explore\_002\_lgbm\_lambdarank\_selected\_forward\_hl30\_train2025.pkl}；004
/ 029 因 GPU 排队 + dataset\_provenance
schema
临时变化没能及时跑完。\textbf{最终上线
6/1 用的是
\texttt{20260601\_explore\_002\_targets.csv}，不是
ensemble}。这是个工程教训：把 ensemble
当成 nice-to-have，没把它包进每日
pipeline。后续 6/2 \textasciitilde{}
6/12 沿用单模型
explore\_002，是诚实的事实而不是设计选择。

\hypertarget{ux4e3bux7ebf-3alpha-stage-ux663eux5f0fux56e0ux5b50ux5e93ux672a-promoteux4f46ux7ea6ux675fux4e86ux8fb9ux754c}{%
\subsubsection{4.3 主线 3：alpha-stage
显式因子库（未
promote，但约束了边界）}\label{ux4e3bux7ebf-3alpha-stage-ux663eux5f0fux56e0ux5b50ux5e93ux672a-promoteux4f46ux7ea6ux675fux4e86ux8fb9ux754c}}

并行另开一条「显式因子库 + agent
pipeline」线（\texttt{alpha-stage/}），目的是补充树模型看不见的可解释
alpha：

\begin{itemize}
\item
  \textbf{Cycle
  1（A001\textasciitilde A006）}：覆盖反转
  / 动量 / 波动率 / 流动性 / 换手率 /
  市值 / 资金流。2026 年内短窗 sanity
  出现「假阳性」：A003 20 日动量 H5
  RankIC 0.10、A005 Amihud 流动性 H5
  RankIC 0.029；扩到 2025-2026 全窗后
  RankIC 反向（-0.052 / -0.029）。Codex
  二审打 1/10，A003 / A005 全部 kill。
\item
  \textbf{回测协议修复}：信号 universe
  与次日成交分离、H\textgreater1
  年化口径修正、long-short
  标记为诊断而非可交易、退出 fillability
  拦掉跌停被卡住的退出。
\item
  \textbf{Cycle
  2（A019\textasciitilde A030）}：跳短期反转、转中期动量
  + 低波动 + 流动性 / 价值、Amihud
  一阶差分等 repair 候选。最强候选
  \textbf{A029（中期动量 + 短回撤 +
  低冲击）} 表达式
  \texttt{rank(rank(ret\_20\_skip5)*(1-rank(ret\_5))*rank(liq\_inv)*rank(low\_vol)*(1-rank(turn)))}。验证集表现：

  \begin{longtable}[]{@{}
    >{\raggedright\arraybackslash}p{(\columnwidth - 10\tabcolsep) * \real{0.17}}
    >{\raggedright\arraybackslash}p{(\columnwidth - 10\tabcolsep) * \real{0.17}}
    >{\raggedright\arraybackslash}p{(\columnwidth - 10\tabcolsep) * \real{0.17}}
    >{\raggedright\arraybackslash}p{(\columnwidth - 10\tabcolsep) * \real{0.17}}
    >{\raggedright\arraybackslash}p{(\columnwidth - 10\tabcolsep) * \real{0.17}}
    >{\raggedright\arraybackslash}p{(\columnwidth - 10\tabcolsep) * \real{0.17}}@{}}
  \toprule
  成本 & val RankIC & 长短年化 & 长短
  Sharpe & long-only 年化 & long-only
  Sharpe \\
  \midrule
  \endhead
  5 bps & 0.045 & +4.92\% & 0.246 &
  +18.06\% & 1.244 \\
  10 bps & 0.045 & +3.50\% & 0.175 &
  +17.33\% & 1.195 \\
  20 bps & 0.045 & +0.65\% & 0.032 &
  +15.89\% & 1.095 \\
  30 bps & 0.045 & -2.20\% & -0.110 &
  +14.45\% & 0.995 \\
  \bottomrule
  \end{longtable}

  长短组合在 30bps 下转负、Sharpe
  跨成本不稳定，未通过 critic gate；但
  long-only top-decile 在
  5\textasciitilde30bps 全档都正、Sharpe
  1.0\textasciitilde1.24（注意：long-only
  是 \textasciitilde330 只 ≠ 比赛要求的
  10 只组合，仅作 spread
  参考）。结论：A029 状态为
  \texttt{repair}，不进入 promote。
\item
  \textbf{GPU 路径}：所有 alpha-stage
  全量回测一律
  \texttt{bash\ scripts/submit\_alpha\_gpu\_backtest.sh}
  走 Slurm，登录节点不直接调 CUDA。
\end{itemize}

这条线没有产出 promote 的
alpha，但它\textbf{约束了主线树模型可以宣称的功能}：报告里不把
long-short 利差当成可交易策略，仅作
spread 诊断。

\hypertarget{ux4e0eux57faux51c6ux5bf9ux6bd4ux9650ux5b9aux53e3ux5f84}{%
\subsubsection{4.4
与基准对比（限定口径）}\label{ux4e0eux57faux51c6ux5bf9ux6bd4ux9650ux5b9aux53e3ux5f84}}

\begin{itemize}
\tightlist
\item
  沪深 300（000300.SH）2025-01
  \textasciitilde{} 2026-04
  期间累计收益约 +9\% / Sharpe
  \textasciitilde0.4，最大回撤 -15\%。
\item
  May 2026 单月：实际上线的 explore\_002
  在 16 个交易日上 Top-10 利差 +177
  bp/天，同期沪深 300 -3.1\%；BT 年化
  +144\% --- 注意：BT 样本只有 15
  个交易日，年化口径要保留谨慎性。
\item
  \textbf{不外推}：长验证集 IC 0.099 是
  LSTM 的成绩，不能用来声称「树模型
  ensemble 在长验证集上 IC
  提升一个数量级」。我们只能说：在 short
  val（May 16 天）这个有限窗口上，从
  LSTM 的「-88.9\% 年化」回到
  explore\_002 的「+144\% 年化」，是
  short val 这一窗口下的差距，长 val
  上没有同口径的树模型 BT 可比。
\end{itemize}

\hypertarget{ux9632ux6cc4ux9732ux4e0eux7a33ux5065ux6027ux81eaux68c0}{%
\subsubsection{4.5
防泄露与稳健性自检}\label{ux9632ux6cc4ux9732ux4e0eux7a33ux5065ux6027ux81eaux68c0}}

\begin{itemize}
\tightlist
\item
  \texttt{tests/test\_no\_leakage.py}
  三项断言：特征严格因果、标签为次日 log
  收益、窗口 z-score 均值 \textless{}
  1e-5。
\item
  \texttt{tests/test\_training\_pipeline.py}
  跑一次 mini 训练 +
  推理，保证训练入口不会因 schema
  偏移而无声降级。
\item
  我们在 v1
  特征里曾踩过一次坑：\texttt{\_cross\_section\_ranks}
  误用未 shift 的 \texttt{ret\_1}，导致
  IC 一度到 0.21（不正常）。修复后回到
  0.04x，并补了对应单测。
\item
  「单一指标陷阱」证据链：explore\_038（Sharpe
  3.09 / IC
  -0.015）、carhart\_ff\_plus\_momentum（IC
  0.044 / 年化
  -21\%）、sklearn\_sgd（RankIC 0.049 /
  年化 -47\%）、wq\_03 vs wq\_06（同
  IC、同向 RankIC，BT 差
  80\%）、rescue\_2025\_rank\_alpha（6
  天 Sharpe 8.3 / RankIC 反向）、event
  reversal（AUC 0.626 / top-k
  期望负）。\textbf{永远要
  IC、RankIC、ICIR、Top-K
  利差、年化、Sharpe、最大回撤六项联合判断}。
\end{itemize}

\hypertarget{ux7ec4ux5408ux56deux6d4bux534fux8baeux4ee3ux7801srcbacktest.py}{%
\subsubsection{\texorpdfstring{4.6
组合回测协议（代码：\texttt{src/backtest.py}）}{4.6 组合回测协议（代码：src/backtest.py）}}\label{ux7ec4ux5408ux56deux6d4bux534fux8baeux4ee3ux7801srcbacktest.py}}

\begin{itemize}
\tightlist
\item
  初始 100 万，n=10 等权建仓，每日换
  k=2，双边手续费 0.0003。
\item
  涨停过滤：估计
  \textbar pct\_chg\textbar{} ≥ 9.5\%
  的不可买。
\item
  基准对比：000300.SH（沪深 300）。
\item
  输出：年化、夏普、最大回撤、净值曲线（\texttt{reports/figures/nav\_compare.png}）。
\item
  \textbf{已知局限}：原版
  \texttt{src/backtest.py} 用次日
  pct\_chg
  做交易过滤，存在轻微未来信息成分；alpha-stage
  的回测脚本已修复，但作业版回测保留这一已知问题并在报告中说明，不在本次大作业中重训整个流水线。
\end{itemize}

\hypertarget{ux6a21ux62dfux4ea4ux6613ux5f85ux586b06-12-ux540eux8865ux5b8c}{%
\subsection{5. 模拟交易（待填，06-12
后补完）}\label{ux6a21ux62dfux4ea4ux6613ux5f85ux586b06-12-ux540eux8865ux5b8c}}

\hypertarget{ux6bcfux65e5ux6d41ux7a0b}{%
\subsubsection{5.1
每日流程}\label{ux6bcfux65e5ux6d41ux7a0b}}

每日 9:00 前运行（实际上线模型为
\texttt{explore\_002} LightGBM
LambdaRank，不是 §3.1 的
transformer；旧版
\texttt{src.predict\_daily\ -\/-model\ transformer}
命令保留但未在比赛中使用）：

\begin{Shaded}
\begin{Highlighting}[]
\CommentTok{\# 1. 同步并构建截至 T 的特征面板}
\ExtensionTok{python} \AttributeTok{{-}m}\NormalTok{ src.data\_loader build{-}panel }\AttributeTok{{-}{-}start}\NormalTok{ 2019{-}01{-}01 }\AttributeTok{{-}{-}end}\NormalTok{ YYYY{-}MM{-}DD}

\CommentTok{\# 2. 用 explore\_002 跑次日 Top{-}N 名单}
\ExtensionTok{python}\NormalTok{ scripts/predict\_strategy\_targets.py }\DataTypeTok{\textbackslash{}}
  \AttributeTok{{-}{-}asof}\NormalTok{ YYYY{-}MM{-}DD }\DataTypeTok{\textbackslash{}}
  \AttributeTok{{-}{-}checkpoint}\NormalTok{ checkpoints/explore\_002\_lgbm\_lambdarank\_selected\_forward\_hl30\_train2025.pkl }\DataTypeTok{\textbackslash{}}
  \AttributeTok{{-}{-}feature{-}set}\NormalTok{ selected }\AttributeTok{{-}{-}direction}\NormalTok{ forward }\AttributeTok{{-}{-}n}\NormalTok{ 10 }\DataTypeTok{\textbackslash{}}
  \AttributeTok{{-}{-}out}\NormalTok{ reports/daily\_logs/YYYYMMDD\_explore\_002\_targets.csv}
\end{Highlighting}
\end{Shaded}

对比前一日 \texttt{*\_targets.csv}
得到买卖清单，限价单挂在前日
\texttt{vwap\ ±\ 0.5\%}，盘中 11:00 /
14:30 各巡查一次，收盘后写
\texttt{reports/daily\_logs/YYYYMMDD\_fills.csv}。

\hypertarget{ux65e5ux5e38ux8bb0ux5f55ux5360ux4f4d}{%
\subsubsection{5.2
日常记录（占位）}\label{ux65e5ux5e38ux8bb0ux5f55ux5360ux4f4d}}

\begin{longtable}[]{@{}llllll@{}}
\toprule
日期 & 目标收益（模型） &
实盘收益（同花顺） & 调仓数 & 未成交 &
备注 \\
\midrule
\endhead
2026-06-01 & & & & & \\
\ldots{} & & & & & \\
\bottomrule
\end{longtable}

\hypertarget{ux6536ux76caux5bf9ux6bd4ux5360ux4f4d}{%
\subsubsection{5.3
收益对比（占位）}\label{ux6536ux76caux5bf9ux6bd4ux5360ux4f4d}}

附同花顺截图：收益曲线 / 持仓 /
调仓记录。

\hypertarget{ux5b9eux9a8cux4eaeux70b9}{%
\subsection{6.
实验亮点}\label{ux5b9eux9a8cux4eaeux70b9}}

\begin{enumerate}
\def\labelenumi{\arabic{enumi}.}
\tightlist
\item
  \textbf{完整的「失败-诊断-换路」过程}：MLP/Transformer→LSTM→LSTM
  在短期失败→LightGBM
  huber→LambdaRank→43
  候选网格→ensemble，每一步换路都有具体证据（IC、Sharpe、Top-K
  bp），不是凭感觉。
\item
  \textbf{IC loss + 日批
  sampler}：训练目标与选股目标对齐，比
  MSE 更适合横截面打分；MLP+MSE IC=0.014
  → LSTM+IC=0.099 是最直接的证据。
\item
  \textbf{窗口内
  z-score（float64）}：规避作业明确禁止的「全量标准化」，并用单测固定行为。
\item
  \textbf{43
  候选系统化网格}：一次性扫完模型/损失/特征/半衰期/训练起点
  5 个维度，发现「2024
  数据有毒、LambdaRank
  显著优于回归、explore\_038 高 Sharpe +
  负 IC 是过拟合」三类非平凡结论。
\item
  \textbf{Ensemble
  多样性}：故意挑两个框架（LGBM/XGB）、两种损失（LambdaRank/pseudoHuber）、两种特征集（selected/alpha\_only），靠
  cross-section rank 平均，而不是 logit
  平均。
\item
  \textbf{防泄露三件套 + CI}：因果断言 +
  标签断言 +
  窗口标准化断言，在历史上确实抓到过一次未
  shift 的
  \texttt{\_cross\_section\_ranks}
  泄露。
\item
  \textbf{盘前 CLI + SKILL.md}（对标
  LLMQuant Awesome Trading
  Agents）：一条命令走完「读最新数据 →
  生成 Top-10 目标 → 输出限价参考 → diff
  买卖清单」。
\item
  \textbf{GPU 路径强制走
  Slurm}：alpha-stage 全量回测一律
  \texttt{bash\ scripts/submit\_alpha\_gpu\_backtest.sh}，登录节点不直接调
  CUDA，避免「在登录节点偷算」的合规风险。
\end{enumerate}

\hypertarget{ux611fux609fux4e0eux601dux8003}{%
\subsection{7.
感悟与思考}\label{ux611fux609fux4e0eux601dux8003}}

\begin{itemize}
\tightlist
\item
  A 股日频 IC 典型 0.02
  \textasciitilde{} 0.08，若验证 IC
  \textgreater{} 0.15 基本可认定存在泄露
  ------ 我们在 v1
  特征中遇到过一次（后发现
  \texttt{\_cross\_section\_ranks}
  误用未 shift 的
  \texttt{ret\_1}，修复后回到 0.04x）。
\item
  模型不是主要瓶颈；\textbf{特征工程 +
  无泄露 + 样本选择}才是。同样的
  LambdaRank 配置仅靠把训练起点从 2024
  换成 2025，Sharpe 就从 -3.06 跳到
  +2.38。
\item
  「长验证集 IC
  高」不等于「比赛期能用」：我们最早信得最多的
  \texttt{lstm\_ic\_large} 在 300 天验证
  IC=0.099，但在 5 月 16 个交易日上变成
  -88\% 年化。最终入选的 LightGBM
  ensemble 是被「短窗口 + regime
  一致性」筛出来的，不是被长验证集 IC
  选出来的。
\item
  单一指标不可信：explore\_038 Sharpe
  3.09 但 IC\textless0；wq\_06 与 wq\_03
  RankIC 都是 0.041 但回测年化差 80
  个百分点。我们最终用
  IC、RankIC、ICIR、Top-K
  利差、年化、Sharpe、最大回撤六项联合判断。
\item
  LLM Agent（外部资料中铺天盖地的多
  Agent 辩论）在选股信号层面远不如一个小
  LSTM/树模型 +
  好特征。但在「每日盘前复盘报告」这类叙事性任务上，Claude
  Code + SKILL.md
  组合体验非常好，本次作业的
  \texttt{skills/a-share-daily-report/SKILL.md}
  即为本地化实践。
\end{itemize}

\hypertarget{ux5206ux5de5}{%
\subsection{8.
分工}\label{ux5206ux5de5}}

\begin{longtable}[]{@{}lll@{}}
\toprule
成员 & 学号 & 分工 \\
\midrule
\endhead
{[}姓名1{]} & {[}学号1{]} &
数据/特征/标签、无泄露单测、报告§2、日志 \\
{[}姓名2{]} & {[}学号2{]} &
模型/训练/回测、SKILL.md、报告§3-4、日志 \\
\bottomrule
\end{longtable}

\hypertarget{ux9644ux5f55-a-ux590dux73b0ux547dux4ee4}{%
\subsection{附录 A
复现命令}\label{ux9644ux5f55-a-ux590dux73b0ux547dux4ee4}}

\begin{Shaded}
\begin{Highlighting}[]
\ExtensionTok{pip}\NormalTok{ install }\AttributeTok{{-}r}\NormalTok{ requirements.txt}

\CommentTok{\# 1. 构建特征 parquet（首次 \textasciitilde{}5{-}10 分钟）}
\ExtensionTok{python} \AttributeTok{{-}m}\NormalTok{ src.data\_loader build{-}panel }\AttributeTok{{-}{-}start}\NormalTok{ 2019{-}01{-}01 }\AttributeTok{{-}{-}end}\NormalTok{ 2026{-}04{-}30}

\CommentTok{\# 2. 深度学习基线 + IC loss 主线}
\ExtensionTok{python} \AttributeTok{{-}m}\NormalTok{ src.train }\AttributeTok{{-}{-}model}\NormalTok{ mlp         }\AttributeTok{{-}{-}loss}\NormalTok{ mse }\AttributeTok{{-}{-}tag}\NormalTok{ mlp\_mse}
\ExtensionTok{python} \AttributeTok{{-}m}\NormalTok{ src.train }\AttributeTok{{-}{-}model}\NormalTok{ lstm        }\AttributeTok{{-}{-}loss}\NormalTok{ ic  }\AttributeTok{{-}{-}tag}\NormalTok{ lstm\_ic\_large }\DataTypeTok{\textbackslash{}}
    \AttributeTok{{-}{-}hidden}\NormalTok{ 256 }\AttributeTok{{-}{-}layers}\NormalTok{ 3 }\AttributeTok{{-}{-}dropout}\NormalTok{ 0.2 }\AttributeTok{{-}{-}epochs}\NormalTok{ 8 }\AttributeTok{{-}{-}batch{-}cap}\NormalTok{ 8192 }\AttributeTok{{-}{-}amp}
\ExtensionTok{python} \AttributeTok{{-}m}\NormalTok{ src.train }\AttributeTok{{-}{-}model}\NormalTok{ transformer }\AttributeTok{{-}{-}loss}\NormalTok{ ic  }\AttributeTok{{-}{-}tag}\NormalTok{ transformer\_ic\_large }\DataTypeTok{\textbackslash{}}
    \AttributeTok{{-}{-}d{-}model}\NormalTok{ 256 }\AttributeTok{{-}{-}heads}\NormalTok{ 8 }\AttributeTok{{-}{-}layers}\NormalTok{ 4 }\AttributeTok{{-}{-}dropout}\NormalTok{ 0.2 }\AttributeTok{{-}{-}epochs}\NormalTok{ 8 }\AttributeTok{{-}{-}batch{-}cap}\NormalTok{ 8192 }\AttributeTok{{-}{-}amp}

\CommentTok{\# 3. LightGBM/XGBoost 主线（包含 43 候选网格）}
\ExtensionTok{python}\NormalTok{ scripts/train\_lgbm\_wq\_short.py             }\CommentTok{\# wq\_00 \textasciitilde{} wq\_06 七候选}
\ExtensionTok{python}\NormalTok{ scripts/explore\_models.py                  }\CommentTok{\# 43 候选网格，\textasciitilde{}28 min on A800}

\CommentTok{\# 4. 跨模型对比 + 回测}
\ExtensionTok{python} \AttributeTok{{-}m}\NormalTok{ src.compare }\AttributeTok{{-}{-}tags}\NormalTok{ mlp\_mse lstm\_ic\_large transformer\_ic\_large transformer\_ic\_mid }\DataTypeTok{\textbackslash{}}
\NormalTok{    lstm\_ic\_h384\_l2\_w20 lstm\_ic\_h384\_l3\_w20 lstm\_ic\_h256\_l3\_w40}
\ExtensionTok{python} \AttributeTok{{-}m}\NormalTok{ src.backtest }\AttributeTok{{-}{-}preds}\NormalTok{ checkpoints/lstm\_ic\_large\_val\_preds.parquet }\DataTypeTok{\textbackslash{}}
    \AttributeTok{{-}{-}out}\NormalTok{ reports/backtest\_lstm\_ic\_large}

\CommentTok{\# 5. 比赛期每日盘前出单（实际上线路径）：{-}{-}asof 是数据截止日}
\ExtensionTok{python}\NormalTok{ scripts/predict\_strategy\_targets.py }\DataTypeTok{\textbackslash{}}
  \AttributeTok{{-}{-}asof}\NormalTok{ 2026{-}05{-}29 }\DataTypeTok{\textbackslash{}}
  \AttributeTok{{-}{-}checkpoint}\NormalTok{ checkpoints/explore\_002\_lgbm\_lambdarank\_selected\_forward\_hl30\_train2025.pkl }\DataTypeTok{\textbackslash{}}
  \AttributeTok{{-}{-}feature{-}set}\NormalTok{ selected }\AttributeTok{{-}{-}direction}\NormalTok{ forward }\AttributeTok{{-}{-}n}\NormalTok{ 10 }\DataTypeTok{\textbackslash{}}
  \AttributeTok{{-}{-}out}\NormalTok{ reports/daily\_logs/20260529\_explore\_002\_targets.csv}

\CommentTok{\# 5.b 旧版 transformer CLI（保留可复现，未用于比赛）：}
\CommentTok{\# python {-}m src.predict\_daily {-}{-}date 2026{-}05{-}29 {-}{-}model transformer {-}{-}tag transformer\_ic\_large {-}{-}n 10}

\CommentTok{\# 6. 单测}
\ExtensionTok{pytest}\NormalTok{ tests/}
\end{Highlighting}
\end{Shaded}

\hypertarget{ux9644ux5f55-b-ux53c2ux8003}{%
\subsection{附录 B
参考}\label{ux9644ux5f55-b-ux53c2ux8003}}

\begin{itemize}
\tightlist
\item
  作业说明：\texttt{大作业\ (1).md}
\item
  USTC 云盘：https://pan.ustc.edu.cn/
  群组「深度学习基础-2026」
\item
  数据说明：zip 内 \texttt{README.md}
\item
  外部地图：https://github.com/LLMQuant/awesome-trading-agents
\end{itemize}

\end{document}
